\subsection*{MT-M23. Every nonabelian Galois top field obstructs weak Mertens and faithful positivity}

Let $L/K$ be a nonabelian Galois extension of number fields. The
classical theorem of Hu--Kaneko--Martin--Schildkraut (HKMS),
\emph{Order of zeros of Dedekind zeta functions}, Theorem1.1,
author version2107.03269v2 p.2, proves unconditionally that
$\zeta_L$ has infinitely many nontrivial zeros of multiplicity
at least two. The complete first proof is \S3 p.5; a second
proof occupies \S4 pp.5--6. Their Theorem1.2, p.2 with proof
\S5 pp.6--8, gives infinitely many zeros of multiplicity at
least three whenever $\operatorname{Gal}(L/K)$ has an irreducible
representation of degree at least three. These are cited classical
results; no new multiplicity theorem is claimed here.

Set
\[
 M_L(x)=\sum_{N\mathfrak a\le x}\mu_L(\mathfrak a),\qquad
 I_L(X)=\int_1^X|M_L(x)/x|^2\,dx.
\]
The following consequences hold without either adopted GRH input:
\[
 \limsup_{X\to\infty}\frac{I_L(X)}{(\log X)^3}>0,
 \qquad
 \limsup_{X\to\infty}\frac{I_L(X)}{\log X}=+\infty.
 \tag{MT-M23a}
\]
If the group has an irreducible representation of degree at least
three, the exponent $3$ in the first assertion can be replaced by
$5$. The canonical full-jet quotient of $\zeta_L$ has no faithful
continuous Hilbert completion on which every convolution extends
boundedly and every $C_{h*h^\sharp}$ is a positive operator.
Nevertheless its nuclear trace construction and adapted absolute
basis exist unconditionally by MT-M18--MT-M19, and $M_H$ makes
the traces of these adjoint-square convolutions nonnegative.

\paragraph{Complete proof of the consequences.}
Choose a multiple zero supplied by HKMS, of actual multiplicity
$m\ge2$, or $m\ge3$ in the stronger case. If $\zeta_L$ has any
off-critical zero, the functional equation and conjugation give
a zero with real part greater than $1/2$. The last assertion of
NT-L16 then gives infinite limsup after division by every fixed
nonnegative power of $\log X$; all displayed conclusions follow.
Otherwise all nontrivial zeros are critical, and our chosen zero
$\rho$ satisfies the exact NT-L16 bound
\[
 \limsup_{X\to\infty}\frac{I_L(X)}{(\log X)^{2m-1}}
 \ge
 \frac{1}{(2m-1)((m-1)!)^2
       |\rho\,\zeta_L^{(m)}(\rho)/m!|^2}>0.
\]
When $2m-1$ equals the claimed exponent this is the first assertion.
When it is greater, take an unbounded sequence on which the
displayed ratio is bounded below by a fixed positive constant;
multiplying by the remaining positive power of $\log X$ gives
infinite limsup. The same argument with exponent one proves
the second assertion. This dichotomy uses no unproved zero
location. Finally MT-M20 proves that a faithful positive Hilbert
realization with the stated properties forces every nontrivial
zero to be simple. The actual multiple zeros contradict that
necessary condition. MT-M18--MT-M19 give the separate basis and
trace statements, and their GRH equivalence gives the conditional
trace positivity under $M_H$.

For the programme this is a finite-monodromy Artin-motive
obstruction to discarding nilpotent spectral blocks. It sharpens
the Kummer examples NT-L17 to all nonabelian Galois top fields.
The cote1/SA23 \texttt{geo-H-lidskii} trace route must retain the
difference between positivity of a scalar trace and positivity
of every operator in a faithful full-multiplicity Hilbert model.
HKMS's general multiplicity prediction for representations of
degree at least four remains their Conjecture6.2, not a theorem
used here. A non-Galois field is covered only when it is itself
the top field of a nonabelian Galois extension of some subfield.

